% !TEX root = ../main.tex
\corrigesChapitre{7}{Cohésion, dissolution et extraction}
\corrigesSouples

\begin{blocLivre}{cor-7-1}
\begin{corrige}{1}{La cohésion du sel}
\demarche{Dans un cristal ionique, chaque ion est entouré d'ions de charge opposée qui l'attirent.}
\begin{questions}
\item \res{6} ions chlorure : quatre dans son plan, un devant et un derrière.
\item Deux ions de signes opposés s'\res{\text{attirent}} ; deux ions de même signe se \res{\text{repoussent}}. Les ions les plus proches
étant de signes opposés, l'attraction l'emporte.
\item Le cristal est électriquement neutre : il contient \res{\text{autant d'ions}} $\ce{Na+}$ que d'ions $\ce{Cl-}$, d'où la formule $\ce{NaCl}$.
\item Les interactions entre ions sont bien plus intenses que les ponts hydrogène qui lient les molécules d'eau dans la
glace : il faut une température \res{\text{bien plus élevée}} pour les vaincre.
\end{questions}
\end{corrige}
\end{blocLivre}
\begin{blocLivre}{cor-7-2}
\begin{corrige}{2}{Écrire des équations de dissolution}
\demarche{La concentration d'un ion vaut son nombre stœchiométrique multiplié par la concentration $c$ du soluté apporté.}
\begin{questions}
\item $\ce{KCl(s) -> K+(aq) + Cl^{-}(aq)}$ ; $\ce{Na2SO4(s) -> 2Na+(aq) + SO4^{2-}(aq)}$ ; $\ce{Al2(SO4)3(s) -> 2Al^{3+}(aq) + 3SO4^{2-}(aq)}$ ;
$\ce{K3PO4(s) -> 3K+(aq) + PO4^{3-}(aq)}$ ; $\ce{Mg(NO3)2(s) -> Mg^{2+}(aq) + 2NO3^{-}(aq)}$.
\item En \unit{mol.L^{-1}} : $\ce{K+}$ et $\ce{Cl-}$ : 0,050 ; $\ce{Na+}$ : 0,100 et $\ce{SO4^{2-}}$ : 0,050 ; $\ce{Al^{3+}}$ : 0,100 et $\ce{SO4^{2-}}$ : 0,150 ;
$\ce{K+}$ : 0,150 et $\ce{PO4^{3-}}$ : 0,050 ; $\ce{Mg^{2+}}$ : 0,050 et $\ce{NO3-}$ : 0,100.
\item Dans la solution de \res{\text{phosphate de potassium}} : $[\ce{K+}]=\qty{0,150}{mol.L^{-1}}$.
\item Charges positives : $3\times0{,}100=0{,}300$ ; charges négatives : $2\times0{,}150=0{,}300$ (en \unit{mol.L^{-1}}) : elles se \res{\text{compensent}}.
\end{questions}
\piege{Écrire $[\ce{Cl-}]=c$ pour le chlorure de calcium : chaque unité $\ce{CaCl2}$ libère deux ions chlorure, donc $[\ce{Cl-}]=2c$. On lit
toujours les proportions sur l'équation de dissolution.}
\end{corrige}
\end{blocLivre}
\begin{blocLivre}{cor-7-3}
\begin{corrige}{3}{Qui forme des ponts hydrogène ?}
\demarche{Il faut un atome H lié à O, N ou F d'un côté, et un doublet non liant de O, N ou F de l'autre.}
\begin{questions}
\item Un atome d'hydrogène lié à un atome \res{\text{très électronégatif}} (O, N, F) et un \res{\text{doublet non liant}} d'un atome O, N ou F
voisin.
\item \res{\ce{H2O}}, \res{\ce{NH3}}, \res{\ce{HF}} et \res{\ce{CH3OH}}. Le méthane n'a pas d'hydrogène lié à O, N ou F.
\item Oui : il ne possède pas d'hydrogène lié à l'oxygène, mais son oxygène porte deux doublets non liants qui peuvent
\res{\text{accueillir}} un hydrogène d'une molécule d'eau.
\item $\ce{CH3-O-H}\cdots\ce{OH2}$ : l'hydrogène du groupe $\ce{-OH}$ du méthanol pointe vers un doublet non liant de l'oxygène de l'eau
(ou l'inverse).
\end{questions}
\schema[corrige]{07-ex03-cor}
\end{corrige}
\end{blocLivre}
\begin{blocLivre}{cor-7-4}
\begin{corrige}{4}{Prévoir une solubilité}
\demarche{Les espèces polaires ou ioniques se dissolvent dans l'eau, les espèces apolaires dans le cyclohexane.}
\begin{questions}
\item Le diiode, apolaire, est plus soluble dans le \res{\text{cyclohexane}}.
\item Le chlorure de sodium, ionique, est soluble dans l'\res{\text{eau}}, qui dissocie et hydrate les ions.
\item Le glucose est très soluble dans l'\res{\text{eau}} : ses groupes $\ce{-OH}$ forment des ponts hydrogène avec elle.
\item La vaseline, apolaire, est soluble dans le \res{\text{cyclohexane}}.
\end{questions}
\retiens{Des entités qui interagissent fortement entre elles se mélangent bien : espèces polaires ou ioniques avec un
solvant polaire, espèces apolaires avec un solvant apolaire.}
\end{corrige}
\end{blocLivre}
\begin{blocLivre}{cor-7-5}
\begin{corrige}{5}{Une ou deux phases ?}
\demarche{Deux liquides non miscibles se superposent : le moins dense au-dessus.}
\begin{questions}
\item L'éthanol est polaire et forme des \res{\text{ponts hydrogène}} avec l'eau.
\item Deux phases : le \res{\text{cyclohexane}} ($d=0{,}78$) \res{\text{au-dessus}}, l'eau au-dessous.
\item Le \res{\text{dichlorométhane}} ($d=1{,}33$) est \res{\text{au-dessous}}, l'eau au-dessus.
\item On ajoute quelques gouttes d'eau : la phase dont le volume \res{\text{augmente}} est la phase aqueuse.
\end{questions}
\schema[corrige]{07-ex05-cor}
\end{corrige}
\end{blocLivre}
\begin{blocLivre}{cor-7-6}
\begin{corrige}{6}{Concentrations des ions en Python}
\demarche{Chaque concentration affichée vaut \code{nb * c}, arrondie au millième.}
\begin{questions}
\item \code{NaCl \{'Na+': 0.05, 'Cl-': 0.05\}}, \code{CaCl2 \{'Ca2+': 0.05, 'Cl-': 0.1\}} et \code{Al2(SO4)3 \{'Al3+': 0.1, 'SO42-': 0.15\}}.
\item \code{'K2SO4': \{'K+': 2, 'SO42-': 1\}} et \code{'Na3PO4': \{'Na+': 3, 'PO43-': 1\}}. Affichage :
\code{K2SO4 \{'K+': 0.2, 'SO42-': 0.1\}} et \code{Na3PO4 \{'Na+': 0.3, 'PO43-': 0.1\}}.
\item Pour chaque solide, on calcule la somme des produits « charge $\times$ nombre stœchiométrique $\times\,c$ » sur tous ses ions : elle
doit être \res{\text{nulle}}.
\item Le glucose est une espèce \res{\text{moléculaire}} : il se dissout sans libérer d'ions.
\end{questions}
\end{corrige}
\end{blocLivre}
\begin{blocLivre}{cor-7-7}
\begin{corrige}{7}{Une solution de sulfate de cuivre}
\demarche{Une mole de cristaux libère une mole d'ions $\ce{Cu^{2+}}$ et une mole d'ions $\ce{SO4^{2-}}$ ; l'eau de cristallisation rejoint le solvant.}
\begin{questions}
\item $M=159{,}6+5\times18{,}0=\res{\qty{249,6}{g.mol^{-1}}}$ ; $\res{\ce{CuSO4(s) -> Cu^{2+}(aq) + SO4^{2-}(aq)}}$.
\item $n=\dfrac{2{,}50}{249{,}6}=\qty{1,00e-2}{mol}$, $c=\qty{0,100}{mol.L^{-1}}$ : $\res{[\ce{Cu^{2+}}]=[\ce{SO4^{2-}}]=\qty{0,100}{mol.L^{-1}}}$.
\item $m=0{,}0500\times0{,}2500\times249{,}6\approx\res{\qty{3,12}{g}}$.
\item $t=0{,}100\times63{,}5=\res{\qty{6,35}{g.L^{-1}}}$.
\end{questions}
\end{corrige}
\end{blocLivre}
\begin{blocLivre}{cor-7-8}
\begin{corrige}{8}{Les sels de déneigement}
\demarche{Une unité $\ce{NaCl}$ libère deux ions, une unité $\ce{CaCl2}$ en libère trois.}
\begin{questions}
\item $\ce{NaCl(s) -> Na+(aq) + Cl^{-}(aq)}$ et $\ce{CaCl2(s) -> Ca^{2+}(aq) + 2Cl^{-}(aq)}$.
\item $n=\dfrac{1000}{58{,}5}=\qty{17,1}{mol}$, soit $\res{\qty{34,2}{mol}}$ d'ions.
\item $n=\dfrac{1000}{111{,}1}=\qty{9,00}{mol}$, soit $\res{\qty{27,0}{mol}}$ d'ions.
\item À masse égale, le \res{\text{chlorure de sodium}} libère plus d'ions. Mais le chlorure de calcium reste efficace à des
températures bien plus basses, et sa dissolution dégage de la chaleur, ce qui accélère la fonte.
\end{questions}
\end{corrige}
\end{blocLivre}
\begin{blocLivre}{cor-7-9}
\begin{corrige}{9}{Choisir un solvant pour la caféine}
\demarche{Un solvant d'extraction doit être non miscible avec l'eau et mieux dissoudre l'espèce que l'eau.}
\begin{questions}
\item L'éthanol est \res{\text{miscible}} avec l'eau : il ne formerait pas deux phases ; de plus, la caféine y est moins soluble.
\item Le dichlorométhane est \res{\text{non miscible}} avec l'eau et la caféine y est environ \res{\text{4,5 fois plus soluble}}.
\item La caféine passe dans la phase de \res{\text{dichlorométhane}}, plus dense que l'eau : elle est \res{\text{au-dessous}}.
\item Travailler sous \res{\text{hotte}}, avec gants et lunettes, et récupérer le solvant usagé dans un bidon de récupération.
\end{questions}
\retiens{Trois critères pour un solvant d'extraction : non miscible avec la phase de départ, meilleure solubilité de
l'espèce à extraire, moindre danger.}
\end{corrige}
\end{blocLivre}
\begin{blocLivre}{cor-7-10}
\begin{corrige}{10}{La vinaigrette}
\demarche{Les molécules d'huile sont apolaires ; l'eau du vinaigre est polaire.}
\begin{questions}
\item Les longues chaînes carbonées de l'huile sont \res{\text{apolaires}} : elles n'interagissent pas avec l'eau, qui forme des ponts
hydrogène.
\item L'\res{\text{huile}}, moins dense que le vinaigre, forme la phase supérieure.
\item Les molécules amphiphiles se placent à la surface des gouttelettes d'huile, partie lipophile dans l'huile, partie
hydrophile dans l'eau : elles \res{\text{stabilisent}} l'émulsion.
\item Savons, shampoings, gels douche, lessives, liquides vaisselle, émulsifiants alimentaires comme la lécithine du jaune
d'œuf de la mayonnaise.
\end{questions}
\end{corrige}
\end{blocLivre}
\begin{blocLivre}{cor-7-11}
\begin{corrige}{11}{Le savon et l'eau calcaire}
\demarche{L'ion carboxylate est amphiphile ; s'il précipite avec les ions calcium, il ne lave plus.}
\begin{questions}
\item Partie hydrophile : le groupe \res{\ce{-COO-}} ; partie lipophile : la chaîne \res{\ce{C17H35-}}.
\item Les chaînes lipophiles pénètrent dans la graisse, les têtes restent dans l'eau : des \res{\text{micelles}} se forment autour des
gouttelettes, entraînées au rinçage.
\item De chaque côté : 36 C, 70 H, 4 O, 1 Ca et une charge totale nulle : l'équation est \res{\text{ajustée}}.
\item Les ions stéarate précipités ne sont plus dissous : ils ne peuvent plus former de micelles, et le dépôt grisâtre
salit le linge.
\end{questions}
\end{corrige}
\end{blocLivre}
\begin{blocLivre}{cor-7-12}
\begin{corrige}{12}{Extraire le diiode d'une eau iodée}
\demarche{Le diiode, apolaire, interagit mieux avec un solvant apolaire ; le solvant choisi doit être non miscible avec l'eau.}
\begin{questions}
\item La molécule $\ce{I2}$ est \res{\text{apolaire}} : elle interagit peu avec les molécules d'eau, polaires, et bien avec celles du
cyclohexane, apolaires.
\item L'éthanol dissout très bien le diiode, mais il est \res{\text{miscible}} avec l'eau : aucune séparation ne serait possible. Le
cyclohexane est non miscible et dissout le diiode près de 70 fois mieux que l'eau.
\item En haut, la \res{\text{phase organique}} (cyclohexane, $d=0{,}78$), \res{\text{violette}} ; en bas, la \res{\text{phase aqueuse}}, presque décolorée.
\item $m=0{,}25\times0{,}020=\qty{5,0}{mg}$ de diiode dans \qty{10}{mL} : $\res{\qty{0,50}{g.L^{-1}}}$, bien inférieur à la solubilité de \qty{20}{g.L^{-1}} :
\res{\text{compatible}}.
\end{questions}
\schema[corrige]{07-ampoule}
\end{corrige}
\end{blocLivre}
\begin{blocLivre}{cor-7-13}
\begin{corrige}{13}{Ponts hydrogène et ébullition}
\demarche{Faire bouillir un liquide, c'est séparer ses molécules : plus les interactions sont intenses, plus il faut chauffer.}
\begin{questions}
\item $\ce{NH3}$, $\ce{H2O}$, $\ce{HF}$ et l'éthanol, qui possèdent un hydrogène lié à N, O ou F.
\item Le soufre est peu électronégatif ($2{,}58-2{,}20=0{,}38$) : $\ce{H2S}$ ne forme pas de ponts hydrogène, l'eau oui. L'eau bout
\res{\qty{160}{\degreeCelsius} \text{ plus haut}}.
\item L'éthanol forme des ponts hydrogène grâce à son groupe $\ce{-OH}$ ; le méthoxyméthane, sans hydrogène lié à l'oxygène, ne
le peut pas entre ses molécules : écart de \res{\qty{102}{\degreeCelsius}}.
\item Le méthane est \res{\text{apolaire}} et ses molécules sont petites : seules de faibles interactions de Van der Waals les lient.
\end{questions}
\end{corrige}
\end{blocLivre}
\begin{blocLivre}{cor-7-14}
\begin{corrige}{14}{Laver le linge en eau calcaire}
\demarche{Chaque ion calcium immobilise deux ions stéarate.}
\begin{questions}
\item $[\ce{Ca^{2+}}]=30\times\num{1.0e-4}=\qty{3,0e-3}{mol.L^{-1}}$ et $n=\num{3.0e-3}\times5{,}0=\res{\qty{1,5e-2}{mol}}$.
\item $n(\text{stéarate})=2\times\num{1.5e-2}=\qty{3,0e-2}{mol}$ et $m=\num{3.0e-2}\times306{,}0\approx\res{\qty{9,2}{g}}$.
\item À \num{5}~°f, six fois moins : environ \res{\qty{1,5}{g}}.
\item Ces espèces \res{\text{immobilisent les ions calcium}} avant qu'ils ne précipitent les tensioactifs : la lessive garde son efficacité
avec une eau dure.
\end{questions}
\end{corrige}
\end{blocLivre}
\begin{blocLivre}{cor-7-15}
\begin{corrige}{15}{Une ou deux extractions ?}
\demarche{La masse se répartit entre les phases : $m=c_\text{aq}V_\text{aq}+5\,c_\text{aq}V_\text{org}$.}
\begin{questions}
\item $\res{m=c_\text{aq}\,(V_\text{aq}+5\,V_\text{org})}$.
\item La masse extraite vaut $5\,c_\text{aq}V_\text{org}$ ; en divisant par $m$ : $\res{\dfrac{5\,V_\text{org}}{V_\text{aq}+5\,V_\text{org}}}$.
\item $\dfrac{5\times40}{100+5\times40}=0{,}67$ : $\res{\qty{0,67}{g}}$.
\item Chaque extraction de \qty{20}{mL} retire $\dfrac{100}{200}=\qty{50}{\percent}$ de la caféine restante : il reste $0{,}5\times0{,}5=\qty{25}{\percent}$, soit
$\res{\qty{0,75}{g}}$ extraits. Le \res{\text{protocole B}} est plus efficace : mieux vaut plusieurs petites extractions qu'une grande.
\end{questions}
\retiens{Pour un même volume de solvant, plusieurs extractions successives avec de petits volumes sont plus efficaces
qu'une seule extraction avec tout le volume.}
\end{corrige}
\end{blocLivre}
